\setcounter{section}{8}

  \zwischenueberschrift{Übungsaufgaben}

\inputaufgabe
{}
{

Es seien
\mathkor {} {M} {und} {N} {}
\definitionsverweis {differenzierbare Mannigfaltigkeiten}{}{}
und
\maabbdisp {\varphi} { M } { N
} {}
eine Abbildung. Es sei
\mathl{M= \bigcup_{i \in I} U_i}{} eine
\definitionsverweis {offene Überdeckung}{}{}
von $M$. Zeige, dass $\varphi$ genau dann
\definitionsverweis {differenzierbar}{}{}
ist, wenn alle Einschränkungen
\mathl{\varphi_i= \varphi \vert_{U_i }}{} differenzierbar sind.

}
{} {}

\inputaufgabe
{}
{

Es sei $M$ eine
\definitionsverweis {differenzierbare Mannigfaltigkeit}{}{}
und
\maabbdisp {\alpha} { U } { V
} {}
eine Karte
\zusatzklammer {also
\mathl{U \subseteq M}{} und
\mathl{V \subseteq \R^n}{} offen} {} {.}
Zeige, dass $\alpha$ ein
\definitionsverweis {Diffeomorphismus}{}{}
ist.

}
{} {}

\inputaufgabe
{}
{

Es sei $M$ eine
\definitionsverweis {differenzierbare Mannigfaltigkeit}{}{} und
\maabbdisp {f,g} { M } {\R
} {}
\definitionsverweis {differenzierbare Funktionen}{}{} auf $M$. Beweise die folgenden Aussagen. \aufzaehlungvier{Die Abbildung
\maabbeledisp {f \times g} { M } {\R^2
} { x } {(f(x),g(x))
} {,} ist differenzierbar.
}{ $f+g$ ist
\definitionsverweis {differenzierbar}{}{.}
}{ $f \cdot g$ ist
\definitionsverweis {differenzierbar}{}{.}
}{Wenn $f$ keine Nullstelle besitzt, so ist auch
\mathl{f^{-1}}{} differenzierbar.
}

}
{} {}

\inputaufgabe
{}
{

Es sei
\mathl{S^2 \subseteq \R^3}{} die
\definitionsverweis {Sphäre}{}{.}
Zeige unter Verwendung der
\definitionsverweis {stereographischen Karten}{}{,}
dass die Einschränkungen der Koordinaten
\mathl{x,y,z}{} des Raumes auf die Sphäre
\definitionsverweis {differenzierbare Funktionen}{}{}
sind.

}
{} {}

\inputaufgabe
{}
{

Es seien
\mathkor {} {M} {und} {N} {}
\definitionsverweis {differenzierbare Mannigfaltigkeiten}{}{}
und
\maabbdisp {\varphi} { M } { N
} {}
eine
\definitionsverweis {differenzierbare Abbildung}{}{.}
Zeige, dass dies einen
\definitionsverweis {Ringhomomorphismus}{}{}
\maabbeledisp {\varphi^*} {C^1(N,\R)} {C^1(M,\R)
} { f } {f \circ \varphi
} {,}
induziert.

}
{} {}

\inputaufgabe
{}
{

Zeige, dass zu
\mathl{m \leq n}{} die Einbettung des Unterraumes
\mathl{\R^m}{} in den $\R^n$, die durch
\mathl{(x_1 , \ldots , x_m) \mapsto (x_1 , \ldots , x_m,0 , \ldots , 0)}{} gegeben ist,
\definitionsverweis {beliebig oft differenzierbar}{}{} ist.

}
{} {}

\inputaufgabe
{}
{

Zeige, dass die offene Zylinderoberfläche
\mathl{S^1 \times {]0,1[}}{} zu
\mathl{S^1 \times \R}{,} zur punktierten Ebene
\mathl{\R^2 \setminus \{(0,0)\}}{} und zu
\mathl{S^2 \setminus \{N,S\}}{}
\definitionsverweis {diffeomorph}{}{}
ist.

}
{} {}

Die nächste Aufgabe verwendet folgende Definition.

Eine Funktion
\maabbdisp {f} { {\mathbb K}^n } { {\mathbb K}
} {}
heißt
\definitionswort {homogen vom Grad}{}
$d$, wenn für jeden Punkt
\mathl{P \in {\mathbb K}^n}{} und jedes
\mathl{\lambda \in {\mathbb K}}{} die Beziehung

\mavergleichskettedisp
{\vergleichskette
{f(\lambda P)
}
{ =} { \lambda^d f(P)
}
{ } {
}
{ } {
}
{ } {
}
}
{}{}{}
gilt.

\inputaufgabe
{}
{

Es sei
\maabbdisp {f} {\R^n } {\R
} {}
eine
\definitionsverweis {stetig differenzierbare}{}{} \definitionsverweis {homogene Funktion}{}{,}
die in der
\definitionsverweis {Faser}{}{}
$F$ über
\mathl{a \neq 0}{}
\definitionsverweis {regulär}{}{}
sei. Zeige, dass jede Faser zu
\mathl{b \neq 0}{} eine zu $F$
\definitionsverweis {diffeomorphe}{}{} \definitionsverweis {Mannigfaltigkeit}{}{}
ist.

}
{} {}

\inputaufgabe
{}
{

Es sei
\mathl{]a,b[}{} ein
\definitionsverweis {offenes Intervall}{}{}
und
\maabbdisp {f} {]a,b[} {\R_{+}
} {}
eine
\definitionsverweis {stetig differenzierbare Funktion}{}{.}
Es sei $M$ die Oberfläche des zugehörigen Rotationskörpers. Zeige, dass diese Menge eine zu einem offenen Zylinder diffeomorphe Mannigfaltigkeit ist.

}
{} {}

\inputaufgabe
{}
{

Zeige, dass eine Ellipsoidoberfläche und die Einheitssphäre
$C^\infty$-\definitionsverweis {diffeomorph}{}{} sind.

}
{} {}

\inputaufgabegibtloesung
{}
{

Man gebe ein Beispiel einer
\definitionsverweis {zweidimensionalen}{}{}
\definitionsverweis {zusammenhängenden}{}{} \definitionsverweis {differenzierbaren Mannigfaltigkeit}{}{}
$M$ und einem Punkt

\mavergleichskette
{\vergleichskette
{ P
}
{ \in }{ M
}
{  }{
}
{    }{
}
{   }{
}
}
{}{}{}
derart, dass
\mathkor {} {M} {und} {M \setminus \{P\}} {}
zueinander
\definitionsverweis {diffeomorph}{}{}
sind.

}
{} {}

\inputaufgabe
{}
{

Zeige, dass die
\definitionsverweis {Antipodenabbildung}{}{}
\maabbeledisp {} {S^n } { S^n
} {(x_1 , \ldots , x_{n+1})} {(-x_1 , \ldots , -x_{n+1})
} {,}
ein
\definitionsverweis {fixpunktfreier}{}{} \definitionsverweis {Diffeomorphismus}{}{}
ist, der zu sich selbst invers ist.

}
{} {}

\inputaufgabegibtloesung
{}
{

Zeige, dass eine
\definitionsverweis {differenzierbare Mannigfaltigkeit}{}{}
$M$ der Dimension

\mavergleichskette
{\vergleichskette
{ n
}
{ \geq }{ 1
}
{  }{
}
{    }{
}
{   }{
}
}
{}{}{}
unendlich viele
\definitionsverweis {Diffeomorphismen}{}{}
\maabbdisp {\varphi} { M } { M
} {}
besitzt.

}
{} {}

Die folgenden Aufgaben sollen erläutern, warum man Mannigfaltigkeiten mit offenen Überdeckungen ansetzt.

Es sei
\mathl{\left(  x_n  \right)_{n \in  \N  }}{} eine
\definitionsverweis {Folge}{}{}
in einem
\definitionsverweis {topologischen Raum}{}{}
$X$. Man sagt, dass die Folge gegen

\mavergleichskette
{\vergleichskette
{ x
}
{ \in }{ X
}
{  }{
}
{    }{
}
{   }{
}
}
{}{}{}
\definitionswort {konvergiert}{,} wenn folgende Eigenschaft erfüllt ist.

Zu jeder
\definitionsverweis {offenen Umgebung}{}{}

\mavergleichskette
{\vergleichskette
{ U
}
{ \subseteq }{ X
}
{  }{
}
{    }{
}
{   }{
}
}
{}{}{}
von $x$ gibt es ein

\mavergleichskette
{\vergleichskette
{ n_0
}
{ \in }{ \N
}
{  }{
}
{    }{
}
{   }{
}
}
{}{}{}
derart, dass für alle

\mavergleichskette
{\vergleichskette
{ n
}
{ \geq }{ n_0
}
{  }{
}
{    }{
}
{   }{
}
}
{}{}{}
die Folgenglieder $x_n$ zu $U$ gehören.

In diesem Fall heißt $x$ der \definitionswort {Grenzwert}{} oder der \definitionswort {Limes}{} der Folge. Dafür schreibt man auch

\mavergleichskettedisp
{\vergleichskette
{  \lim_{n \rightarrow \infty}  x_n
}
{ =} { x
}
{ } {
}
{ } {
}
{ } {
}
}
{}{}{.}
Wenn die Folge einen Grenzwert besitzt, so sagt man auch, dass sie \definitionswort {konvergiert}{}
\zusatzklammer {ohne Bezug auf einen Grenzwert} {} {,}
andernfalls, dass sie \definitionswort {divergiert}{.}

\inputaufgabe
{}
{

Es sei $M$ eine
\definitionsverweis {topologische Mannigfaltigkeit}{}{}
und
\mathl{\left(   x_n  \right)_{n \in \N }}{} eine
\definitionsverweis {Folge}{}{}
in $M$. Zeige, dass die Folge genau dann
\definitionsverweis {konvergiert}{}{,}
wenn es ein
\definitionsverweis {Kartengebiet}{}{}
von $M$ gibt, das fast alle Glieder der Folge enthält und so, dass die entsprechende Bildfolge im Kartenbild konvergiert.

}
{} {}

\inputaufgabe
{}
{

Es sei $M$ eine
\definitionsverweis {topologische Mannigfaltigkeit}{}{}
und
\maabbdisp {\alpha_i} {U_i } { V_i
} {}
eine Familie von
\definitionsverweis {Karten}{}{}
mit den
\definitionsverweis {Übergangsabbildungen}{}{}
\maabbdisp {\varphi_{ij} = \alpha_j \circ (\alpha_i)^{-1}} {V_i \cap \alpha_i(U_i \cap U_j) } { V_j \cap \alpha_j(U_i \cap U_j)
} {.}
Zeige, dass man aus der Familie der

\mathbed {V_i} {}
 {i \in I} {}
 {} {} {} {,}
den Teilmengen
\mathl{V_{ij} \subseteq V_i}{} und den Übergangsabbildungen
\maabbdisp {\varphi_{ij}} {V_{ij}} { V_{ji}
} {}
die Mannigfaltigkeit $M$ rekonstruieren kann.

a) Betrachte auf

\mavergleichskettedisp
{\vergleichskette
{N
}
{ \defeq} { \biguplus_{i \in I} V_i
}
{ } {
}
{ } {
}
{ } {
}
}
{}{}{}
die
\definitionsverweis {Äquivalenzrelation}{}{,}
unter der zwei Punkte
\mathl{P\in V_i}{} und
\mathl{Q \in V_j}{} gleich sind, wenn sie unter
\mathl{\varphi_{ij}}{} ineinander abgebildet werden.

b) Versehe die Quotientenmenge
\mathl{N/ \sim}{} mit einer geeigneten Topologie.

c) Definiere auf
\mathl{N/ \sim}{} Karten.

d) Zeige, dass
\mathkor {} {M} {und} {N/ \sim} {}
\definitionsverweis {homöomorph}{}{}
sind.

}
{} {}

Das in der vorstehenden Aufgabe beschriebene Konstruktionsverfahren für eine Mannigfaltigkeit funktioniert für eine Familie von offenen Teilmengen im $\R^n$ mit Übergangsabbildungen, die die Kozykelbedingung aus
[[Mannigfaltigkeit/Übergangsabbildungen/Kozykelbedingung/Aufgabe|Aufgabe 7.11]]
erfüllen. Allerdings ist der dabei entstehende topologische Raum nicht ohne weiteres ein Hausdorff-Raum. Man spricht vom \stichwort {offenen Verkleben} {} von Räumen.

\inputaufgabe
{}
{

Wir betrachten die reelle Gerade zweifach, also

\mavergleichskette
{\vergleichskette
{G_1
}
{ = }{\R
}
{  }{
}
{    }{
}
{   }{
}
}
{}{}{}
und

\mavergleichskette
{\vergleichskette
{G_2
}
{ = }{\R
}
{  }{
}
{    }{
}
{   }{
}
}
{}{}{}
zusammen mit der Verklebungsabbildung
\maabbeledisp {\varphi} {G_1 \setminus \{0\} } { G_2 \setminus \{0\}
} { x } { x^{-1}
} {.}
Es sei $M$ der entstehende topologische Raum gemäß der in
[[Topologische Mannigfaltigkeit/Rekonstruktion aus Karten/Aufgabe|Aufgabe 8.15]]
beschriebenen Konstruktion. Zeige, dass $M$
\definitionsverweis {homöomorph}{}{}
zur
$1$-\definitionsverweis {Sphäre}{}{}
ist.

}
{} {}

\inputaufgabe
{}
{

Wir betrachten die reelle Gerade zweifach, also

\mavergleichskette
{\vergleichskette
{G_1
}
{ = }{\R
}
{  }{
}
{    }{
}
{   }{
}
}
{}{}{}
und

\mavergleichskette
{\vergleichskette
{G_2
}
{ = }{\R
}
{  }{
}
{    }{
}
{   }{
}
}
{}{}{}
zusammen mit der Verklebungsabbildung
\maabbdisp {\operatorname{Id}} {G_1 \setminus \{0\} } { G_2 \setminus \{0\}
} {.}
Es sei $M$ der entstehende topologische Raum gemäß der in
[[Topologische Mannigfaltigkeit/Rekonstruktion aus Karten/Aufgabe|Aufgabe 8.15]]
beschriebenen Konstruktion. Zeige, dass $M$ keine
\definitionsverweis {Mannigfaltigkeit}{}{}
ist.

}
{} {}

  \zwischenueberschrift{Aufgaben zum Abgeben}

In der folgenden Aufgabe interpretiere man ${\mathbb C}$ als $\R^2$.

\inputaufgabe
{4}
{

Wir betrachten die Abbildung
\maabbeledisp {} { {\mathbb C}^2} { {\mathbb C}
} {(z,w)} { zw
} {.}
Für welche Punkte
\mathl{u \in {\mathbb C}}{} ist die Faser über $u$ eine Mannigfaltigkeit? Man gebe jeweils eine möglichst einfache Beschreibung des Diffeomorphietyps.

}
{} {}

\inputaufgabe
{6}
{

Es seien zwei Punkte
\mathkor {} {P} {und} {Q} {}
auf der
\definitionsverweis {Einheitssphäre}{}{}
gegeben. Zeige, dass es einen
\definitionsverweis {Diffeomorphismus}{}{}
der Sphäre in sich gibt, der $P$ in $Q$ überführt.

}
{} {}

\inputaufgabe
{4 (1+1+2)}
{

Es sei $M$ eine
\definitionsverweis {differenzierbare Mannigfaltigkeit}{}{.}
Zu jeder offenen Teilmenge

\mavergleichskette
{\vergleichskette
{ U
}
{ \subseteq }{ M
}
{  }{
}
{    }{
}
{   }{
}
}
{}{}{}
betrachten wir die Menge
\mathl{C^1(U,\R)}{} der
\definitionsverweis {differenzierbaren Funktionen}{}{}
auf $U$. Es sei

\mavergleichskette
{\vergleichskette
{ M
}
{ = }{ \bigcup_{i \in I} U_i
}
{  }{
}
{    }{
}
{   }{
}
}
{}{}{}
eine
\definitionsverweis {offene Überdeckung}{}{.}
\aufzaehlungdreiabc{Zeige, dass zu

\mavergleichskette
{\vergleichskette
{ V
}
{ \subseteq }{ U
}
{  }{
}
{    }{
}
{   }{
}
}
{}{}{}
offen und

\mavergleichskette
{\vergleichskette
{ f
}
{ \in }{ C^1(U,\R)
}
{  }{
}
{    }{
}
{   }{
}
}
{}{}{}
auch die
\definitionsverweis {Einschränkung}{}{}
$f \vert_V$ zu
\mathl{C^1(V,\R)}{} gehört.
}{Es sei

\mavergleichskette
{\vergleichskette
{ f
}
{ \in }{ C^1(M,\R)
}
{  }{
}
{    }{
}
{   }{
}
}
{}{}{.}
Zeige, dass

\mavergleichskette
{\vergleichskette
{ f
}
{ = }{ 0
}
{  }{
}
{    }{
}
{   }{
}
}
{}{}{}
genau dann ist, wenn sämtliche Einschränkungen

\mavergleichskette
{\vergleichskette
{ f \vert_{U_i }
}
{ = }{ 0
}
{  }{
}
{    }{
}
{   }{
}
}
{}{}{}
sind.
}{Es sei eine Familie

\mavergleichskette
{\vergleichskette
{ f_i
}
{ \in }{ C^1(U_i,\R)
}
{  }{
}
{    }{
}
{   }{
}
}
{}{}{}
von Funktionen gegeben, die die \anfuehrung{Verträglichkeitsbedingung}{}

\mavergleichskettedisp
{\vergleichskette
{ f_i \vert_{U_i \cap U_j }
}
{ =} { f_j \vert_{U_i \cap U_j }
}
{ } {
}
{ } {
}
{ } {
}
}
{}{}{}
für alle $i,j$ erfüllen. Zeige, dass es ein

\mavergleichskette
{\vergleichskette
{ f
}
{ \in }{ C^1(M,\R)
}
{  }{
}
{    }{
}
{   }{
}
}
{}{}{}
gibt mit

\mavergleichskette
{\vergleichskette
{ f \vert_{U_i }
}
{ = }{ f_i
}
{  }{
}
{    }{
}
{   }{
}
}
{}{}{}
für alle $i$.
}

}
{} {}

\inputaufgabe
{5}
{

Es sei $M$ eine
\definitionsverweis {differenzierbare Mannigfaltigkeit}{}{,} die mindestens zwei Elemente besitze. Zeige, dass es
\definitionsverweis {differenzierbare Funktionen}{}{}
\maabbdisp {f,g} { M } {\R
} {} gibt mit
\mathl{f,g \neq 0}{,} aber
\mathl{fg=0}{.}

}
{} {}

 [[Kategorie:Latexseite]]
